% ---------------------------------------------------------------
% CHAPTER 3: DOJO TRAINING PHASE
% ---------------------------------------------------------------
\chapter{Training Phase: DOJO Learning}

\section{Overview}

The program commenced with a structured pedagogical curriculum administered by 
seasoned technical mentors. The initiative targeted advancement of idiomatic, 
comprehensively validated GoLang implementations satisfying predefined specifications 
with maximized test case protection. Beyond Go language mastery, the curriculum 
introduced client-side development methodologies using React, foundational cloud 
infrastructure understanding, and emerging machine learning capabilities --- 
delivering comprehensive technical preparation preceding project-based work.

\begin{figure}[H]
\centering
\begin{tikzpicture}[node distance=0.6cm]
  \node[stepblock] (step1) {STEP 1: GoLang Fundamentals};
  \node[stepblock, below=of step1] (step2)
    {STEP 2: Concurrency --- Goroutines \& Channels};
  \node[stepblock, below=of step2] (step3)
    {STEP 3: Gin Web Framework \& RESTful APIs};
  \node[stepblock, below=of step3] (step4)
    {STEP 4: TDD --- Driver-Navigator \& Ping-Pong Pairing};
  \node[stepblock, below=of step4] (step5)
    {STEP 5: Docker --- Containerisation \& Image Management};
  \node[stepblock, below=of step5] (step6)
    {STEP 6: React --- Frontend Development};
  \node[stepblock, below=of step6] (step7)
    {STEP 7: Cloud Computing Fundamentals};
  \node[stepblock, below=of step7] (step8)
    {STEP 8: AI/ML --- Transformers, RAG \& Fine-Tuning};
  \node[widegreenblock, below=of step8] (done)
    {OUTPUT: Full-Stack Engineering Foundation with 80\%+ Test Coverage};
  \draw[arrow] (step1) -- (step2);
  \draw[arrow] (step2) -- (step3);
  \draw[arrow] (step3) -- (step4);
  \draw[arrow] (step4) -- (step5);
  \draw[arrow] (step5) -- (step6);
  \draw[arrow] (step6) -- (step7);
  \draw[arrow] (step7) -- (step8);
  \draw[arrow] (step8) -- (done);
\end{tikzpicture}
\caption{DOJO Learning Phase --- Progression Flow}
\end{figure}

\section{Key Learning Areas}

\subsection{GoLang Fundamentals}

Fundamental GoLang principles studied encompassed primitive data representations, 
program flow direction, type systems, concurrent execution units, communication structures, 
and idiomatic exception management utilizing the \texttt{error} interface. Emphasis 
was positioned on authoring Go source utilizing official language specifications 
and established idioms as published in the Effective-Go publication.

\subsection{Concurrency Patterns}

GoLang's concurrent processing architecture --- established on Communicating Sequential 
Processes conventions --- received instruction via pragmatic demonstrations with concurrent 
execution units and information exchange conduits. This theoretical knowledge demonstrated 
importance when constructing competing reservation processing in SkyFox Cinemas and 
establishing the multi-subscriber recovery solution.

\subsection{Gin Web Framework}

The Gin HTTP foundation was employed to produce economical application programming interfaces. 
Navigation definition, middleware arrangement, JSON payload handling, information validation, 
and formalized error management design inside Gin received comprehensive examination.

\subsection{Test-Driven Development (TDD)}
\begin{figure}[H]
\centering
\begin{tikzpicture}[node distance=0.5cm]

  % ---- Left column: Driver-Navigator ----
  \node[rectangle, fill=idfcblue, text=white, minimum width=5.5cm,
        minimum height=0.7cm, font=\bfseries\small, rounded corners=3pt]
    (dn_title) {Driver-Navigator Pairing};
  \node[block, below=0.4cm of dn_title, text width=5.2cm] (dn1)
    {Navigator designs the approach \& guides};
  \node[block, below=0.4cm of dn1, text width=5.2cm] (dn2)
    {Driver implements the code};
  \node[yellowblock, below=0.4cm of dn2, text width=5.2cm] (dn3)
    {Roles swap every 10--15 minutes};

  % ---- Right column: Ping-Pong ----
  \node[rectangle, fill=idfcblue, text=white, minimum width=5.5cm,
        minimum height=0.7cm, font=\bfseries\small, rounded corners=3pt,
        right=1.0cm of dn_title]
    (pp_title) {Ping-Pong Pairing};
  \node[block, below=0.4cm of pp_title, text width=5.2cm] (pp1)
    {Person A writes a \textbf{failing} test (Red)};
  \node[greenblock, below=0.4cm of pp1, text width=5.2cm] (pp2)
    {Person B makes the test \textbf{pass} (Green)};
  \node[yellowblock, below=0.4cm of pp2, text width=5.2cm] (pp3)
    {Refactor $\rightarrow$ swap roles $\rightarrow$ repeat};

  % ---- Arrows ----
  \draw[arrow] (dn1) -- (dn2);
  \draw[arrow] (dn2) -- (dn3);
  \draw[arrow] (pp1) -- (pp2);
  \draw[arrow] (pp2) -- (pp3);

\end{tikzpicture}
\caption{TDD Methodologies: Driver-Navigator vs.\ Ping-Pong Pairing}
\end{figure}
\subsection{Docker Basics}

An overview of containerization employing Docker technology was furnished, examining 
image construction specifications, hierarchical storage structures, operational lifecycle 
administration, information persistence approaches, and variable injection mechanisms --- 
establishing readiness for downstream deployment automation work.

\subsection{React and Frontend Development}

Frontend implementation methodologies received instruction utilizing React framework, 
analyzing component-driven structural patterns, information handling via function-based 
storage, and consumption of data transfer protocols originating from Go service providers. 
This foundation immediately supported the SkyFox Cinemas program, facilitating development 
of patron engagement environments and administrative oversight dashboard using React.

\subsection{Cloud Computing}

An introduction to distributed cloud delivery infrastructure was presented, including 
provisioning approaches (Infrastructure services, Platform services, Application services), 
packaged execution environments, and cloud-compatible solution topology. This discussion 
contextualized how commercial financial platforms are provisioned and expanded at IDFC 
FIRST Bank.

\subsection{AI and Machine Learning Concepts}

The DOJO similarly addressed contemporary artificial intelligence methodologies relevant 
to software production environments. Subject matter contained the Attention architecture 
fundamental to generative language systems, Documentation-Informed Content Production 
for constraining outputs with institutional information, and parameter fine-tuning for 
application-specific implementations. This education developed consciousness regarding 
the incorporation of AI capabilities into contemporary monetary and technology sector 
manufacturing platforms.

\section{DOJO Outcomes}

\begin{figure}[H]
\centering
\begin{tikzpicture}[node distance=0.45cm]
  \node[rectangle, fill=idfcblue, text=white, minimum width=12.5cm,
        minimum height=0.8cm, font=\bfseries, rounded corners=4pt]
    (title) {Skills Acquired and Their Application};
  \node[wideblock, below=0.5cm of title, text width=12.2cm,
        minimum height=1cm] (o1)
    {\textbf{GoLang + Gin + TDD:} Applied directly in SkyFox Cinemas backend
     and the Resilient Event Consumer Framework};
  \node[wideblock, below=0.4cm of o1, text width=12.2cm,
        minimum height=1cm] (o2)
    {\textbf{Concurrency:} Goroutine patterns used in concurrent seat booking
     and multi-consumer event processing};
  \node[wideblock, below=0.4cm of o2, text width=12.2cm,
        minimum height=1cm] (o3)
    {\textbf{React:} Used to build the SkyFox Cinemas customer and
     administrator UI across three sprints};
  \node[wideblock, below=0.4cm of o3, text width=12.2cm,
        minimum height=1cm] (o4)
    {\textbf{Docker + Cloud:} Informed CI/CD pipeline design and
     containerised service deployment on goCD};
  \node[widegreenblock, below=0.4cm of o4, text width=12.2cm,
        minimum height=1cm] (o5)
    {\textbf{AI/ML:} Provided foundational awareness of Transformer, RAG,
     and fine-tuning patterns used in enterprise AI integration};
\end{tikzpicture}
\caption{DOJO Training --- Skills and Their Application}
\end{figure}